\chapter{Client Interaction Log}
\label{app:client-log}

Cột Ngày ghi thời điểm client trả lời. Email vòng 2 (08/10) gồm 10 câu: Q10--Q17, Q19, Q20.

\kiemtra{Các dòng ngày 08/10 lấy từ file \texttt{Phản hồi 8\_10.txt}; gửi lúc 16:39 bởi Nguyễn Đức Nghĩa, thầy trả lời lúc 19:58 (do PM cung cấp). Ảnh bằng chứng \texttt{Q10-Q20-2026-10-08-email-tra-loi.png} lưu trên Drive của nhóm (Phụ lục~\ref{app:repository}). Cột \enquote{Tác động} của vòng này do AI đề xuất, TV2 rà lại.}

\kiemtra{Email vòng 2 (08/10) mở đầu bằng câu cảm ơn thầy \enquote{đã nhận xét nhóm ở buổi Milestone 1}. Theo biên bản CP (tab Biên bản CP, dòng M1), buổi 05/10 thầy không gọi nhóm 9 trình bày nên không có nhận xét M1; câu mở đầu trên là nguyên văn email đã gửi, không phải bằng chứng có nhận xét M1 của thầy. TV1 xác nhận cách ghi chú này.}

\kiemtra{Dòng Q21, Q22: câu hỏi theo email TV1 gửi 23/10 15:10, trả lời theo ảnh chụp email TV1 cung cấp 27/10; TV1 lưu ảnh lên Drive (\texttt{Q21-Q22-2026-10-25-email-tra-loi.png}) và dán hai dòng này vào tab Client Log.}

\todo{Bổ sung các vòng trao đổi tiếp theo (trao đổi tại Progress Check 09/11; change event).}

\begingroup\tblsetup\tblzebra
\begin{longtable}{L{2.0cm}L{3.4cm}L{4.8cm}L{3.8cm}}
\caption{Client Interaction Log}\label{tab:client-log}\\
\toprule
\rowcolor{tblheadbg}\textbf{Ngày} & \textbf{Câu hỏi của nhóm} & \textbf{Trả lời của client} & \textbf{Tác động tới dự án} \\
\midrule
\endfirsthead
\toprule
\rowcolor{tblheadbg}\textbf{Ngày} & \textbf{Câu hỏi của nhóm} & \textbf{Trả lời của client} & \textbf{Tác động tới dự án} \\
\midrule
\endhead
\bottomrule
\endlastfoot
24/09/2026 07:48 & \textbf{P01.} Ai đóng vai client/sponsor? Ai duyệt scope và thay đổi? & Trong phạm vi BTL, thầy đóng vai client. Nội dung cần làm rõ về scope, priority hoặc thay đổi yêu cầu thì trao đổi trực tiếp với thầy. & Project authority: người quyết định scope, mức ưu tiên và thay đổi yêu cầu là thầy Nguyễn Minh Tâm. \\
24/09/2026 07:48 & \textbf{P02.} Liên hệ client qua kênh nào? Giới hạn và thời gian phản hồi? & Trao đổi qua email. Không giới hạn số câu hỏi nhưng nên tổng hợp trước khi gửi, tránh nhiều email lẻ. Không cam kết thời gian phản hồi cố định; nhóm cần chủ động liên hệ sớm. & Mỗi vòng hỏi gom câu hỏi vào một email. Vòng 2 hỏi qua email sau buổi Milestone 1. \\
24/09/2026 07:48 & \textbf{P03.} Có nên khảo sát sinh viên thật? Trình bày tách khỏi Client Log? & Không. & Không khảo sát sinh viên. Stakeholder analysis dựa trên trao đổi với client và phân tích của nhóm. \\
24/09/2026 07:48 & \textbf{P04.} Hạn nộp và định dạng Milestone 1? & Đã có thông báo trên lớp. & Theo thông báo trên lớp: Milestone 1 là buổi trình bày trên lớp thứ Hai 05/10/2026. Nhóm phát hiện ngày 02/10 (issue I01). \\
24/09/2026 07:48 & \textbf{P05.} Có liên hệ client trong tuần nghỉ để duyệt baseline không? & Trong tuần nghỉ vẫn liên hệ qua email được. Nhóm không cần chờ client "duyệt" toàn bộ Baseline V1.0; baseline là sản phẩm lập kế hoạch của nhóm dựa trên thông tin, assumptions và constraints đã thu thập. Điểm nào ảnh hưởng đáng kể tới scope hoặc planning thì hỏi trước khi chốt baseline. & Nhóm tự chốt Baseline V1.0. Trước khi chốt, gửi email hỏi thầy các điểm ảnh hưởng đáng kể tới scope hoặc planning. \\
24/09/2026 07:48 & \textbf{P06.} Bằng chứng cho Client Log cần dạng nào? & Không yêu cầu biên bản có chữ ký/xác nhận riêng. Nhóm cần ghi lại được thời điểm trao đổi, câu hỏi, phản hồi của client. & Bằng chứng Client Log là ảnh chụp email; cột Ngày ghi kèm giờ. \\
02/10/2026 19:48 & \textbf{Q08.} Tiêu chí nghiệm thu là gì? Thiếu thời gian thì ưu tiên gì? & Nghiệm thu ưu tiên các workflow cốt lõi: tạo/quản lý sự kiện $\rightarrow$ sinh viên đăng ký/hủy đăng ký $\rightarrow$ quản lý số lượng và danh sách chờ $\rightarrow$ điểm danh $\rightarrow$ xem kết quả/báo cáo. Thông báo/nhắc lịch là chức năng bổ trợ; nếu chưa hoàn thiện đầy đủ thì không ảnh hưởng nhiều bằng các chức năng cốt lõi trên. & Xác định thứ tự ưu tiên scope: 5 workflow cốt lõi là Must; thông báo/nhắc lịch ưu tiên thấp hơn. \\
02/10/2026 19:48 & \textbf{Q03.} Điểm danh làm thế nào? Có check-out? Mạng yếu xử lý ra sao? & Có thể sử dụng QR để điểm danh, đồng thời hỗ trợ nhập mã hoặc xác nhận thủ công làm phương án dự phòng. Không bắt buộc check-out. & Điểm danh bằng QR, có phương án nhập mã hoặc xác nhận thủ công. Check-out ngoài scope. \\
02/10/2026 19:48 & \textbf{Q04.} Giới hạn chỗ, danh sách chờ, hạn chót hủy? & Sự kiện có giới hạn số lượng và có waitlist. Khi có chỗ trống, hệ thống tự động chuyển người đầu tiên trong waitlist sang danh sách đăng ký và gửi thông báo, không cần BTC duyệt lại. Nhóm có thể đề xuất deadline hủy đăng ký phù hợp và ghi rõ giả định. & Giới hạn số chỗ và waitlist tự động (kèm thông báo khi được chuyển lên) là yêu cầu. Nhóm cần tự đặt deadline hủy đăng ký và ghi là assumption. \\
02/10/2026 19:48 & \textbf{Q01.} Ai được tạo sự kiện? Có cần duyệt? & Có thể giả định các đơn vị như Khoa, CLB/Đội/Nhóm và các đơn vị/phòng ban của trường có thể tổ chức sự kiện. Sự kiện cần được người có quyền quản lý/duyệt xác nhận trước khi công bố. Nhóm chủ động đề xuất role cụ thể cho workflow này. & Thêm luồng duyệt sự kiện trước khi công bố. Nhóm cần tự đề xuất các role (đơn vị tổ chức, người duyệt\ldots) và ghi là assumption. \\
02/10/2026 19:48 & \textbf{Q02.} Đăng nhập bằng gì? & Có thể giả định sử dụng tài khoản trường. Không yêu cầu tích hợp SSO thực tế; nhóm có thể mô phỏng bằng email trường hoặc dữ liệu tài khoản mẫu. & Tích hợp SSO thật ngoài scope; đăng nhập bằng tài khoản trường mô phỏng (email trường hoặc tài khoản mẫu). \\
02/10/2026 19:48 & \textbf{Q05.} Báo cáo sau sự kiện gồm gì, định dạng nào? & Báo cáo cơ bản gồm danh sách đăng ký, danh sách tham dự/điểm danh và thống kê số lượng người tham gia. Có thể hỗ trợ xuất Excel/CSV. Không yêu cầu tích hợp trực tiếp với hệ thống điểm rèn luyện trong phạm vi hiện tại. & Phạm vi báo cáo: danh sách đăng ký, danh sách điểm danh, thống kê số người; xuất Excel/CSV. Tích hợp điểm rèn luyện ngoài scope. \\
02/10/2026 19:48 & \textbf{Q06.} Quy mô: sự kiện/tháng, người/sự kiện, check-in đồng thời? & Hiện không có số liệu thực tế về số sự kiện/tháng, số người tham gia hoặc tải đồng thời. Nhóm có thể đưa ra giả định hợp lý để phục vụ thiết kế và kiểm thử, nhưng cần ghi rõ là assumption. & Không có số liệu quy mô. Nhóm cần tự đặt giả định (số sự kiện/tháng, số người/sự kiện, tải check-in đồng thời) và ghi là assumption. \\
02/10/2026 19:48 & \textbf{Q09.} Ràng buộc dữ liệu cá nhân, hosting, nền tảng? & Cần có phân quyền và bảo vệ thông tin cá nhân cơ bản của sinh viên. Không có ràng buộc cụ thể về hosting/nền tảng. Web responsive là đủ, không bắt buộc native mobile app. & Phân quyền và bảo vệ thông tin cá nhân cơ bản là yêu cầu. Không ràng buộc hosting. Web responsive; native app ngoài scope. \\
02/10/2026 19:48 & \textbf{Q07.} Hiện dùng công cụ gì? Tổng hợp mất bao lâu? & Có thể giả định quy trình hiện tại sử dụng các công cụ như form/bảng tính và có một số bước tổng hợp thủ công. & Quy trình hiện tại (giả định theo client): form/bảng tính, tổng hợp thủ công. Nhóm cần tự đặt mốc định lượng nếu dùng cho MoE và business case, ghi là assumption. \\
08/10/2026 19:58 & \textbf{Q10.} Phạm vi. Nhóm hiểu phạm vi cốt lõi gồm: tạo/chỉnh sửa sự kiện $\rightarrow$ duyệt trước khi công bố $\rightarrow$ đăng ký/hủy $\rightarrow$ giới hạn chỗ + waitlist + tự động chuyển người đầu waitlist $\rightarrow$ QR check-in (có nhập mã/xác nhận thủ công dự phòng) $\rightarrow$ báo cáo và xuất Excel/CSV. Nhóm không làm: check-out bắt buộc, SSO thật của trường, ứng dụng native Android/iOS, tích hợp điểm rèn luyện, đồng bộ offline đầy đủ. Thầy xác nhận cách hiểu này đã đúng chưa ạ? & Đồng ý & Chốt scope và exclusions làm đầu vào PBS, Product Description và Baseline V1.0. \\
08/10/2026 19:58 & \textbf{Q11.} Mục tiêu và thước đo thành công. Nhóm đang đặt 4 mục tiêu sau: Không cho đăng ký trùng hoặc vượt số chỗ; waitlist chuyển người theo đúng thứ tự vào trước ra trước. Thời gian check-in trung vị $\leq$ 5 giây trong môi trường test của nhóm; có phương án thủ công khi QR không dùng được; chặn check-in trùng. Với 500 lượt đăng ký, tổng hợp điểm danh và xuất Excel/CSV trong $\leq$ 30 giây, số liệu khớp dữ liệu nguồn. Mọi kiểm thử phân quyền đều đạt; sự kiện chưa được duyệt thì không công bố được; không còn lỗi nghiêm trọng về quyền riêng tư/kiểm soát truy cập khi nghiệm thu. Các mốc hiện trạng "15 giây/người" (check-in) và "20 phút/sự kiện" (tổng hợp sau sự kiện) chỉ là giả định lập kế hoạch của nhóm, không phải số liệu thật. Thầy có đồng ý dùng các mục tiêu này làm tiêu chí thành công của dự án không ạ, hay cần điều chỉnh ngưỡng nào? & Đồng ý & Chốt 4 objective và MoE làm tiêu chí thành công; mốc 15 giây/người và 20 phút/sự kiện vẫn là planning assumption. \\
08/10/2026 19:58 & \textbf{Q12.} Role duyệt sự kiện. Nhóm đề xuất dùng một role riêng là Admin/Approver để duyệt sự kiện trước khi công bố, và Organizer không được tự duyệt sự kiện do mình tạo. Đề xuất này có phù hợp với thầy không ạ? & Mình đồng ý với nguyên tắc Organizer không được tự duyệt sự kiện do mình tạo. Tuy nhiên, trong MVP không cần tách thêm một role độc lập nếu không cần thiết. Có thể sử dụng role Admin/System Administrator để thực hiện chức năng duyệt. & Admin duyệt, không thêm role riêng; Organizer không được tự duyệt sự kiện của mình. \\
08/10/2026 19:58 & \textbf{Q13.} Hạn hủy đăng ký. Nhóm đề xuất một mốc cố định cho toàn hệ thống: sinh viên được tự hủy đăng ký đến 24 giờ trước giờ bắt đầu sự kiện. Sau mốc này hệ thống không cho sinh viên tự hủy, nếu có lý do chính đáng thì sinh viên liên hệ Organizer. Thầy thấy cách này có phù hợp không ạ, hay thầy muốn mốc khác hoặc để Organizer cấu hình theo từng sự kiện? & Đồng ý & Hạn hủy cố định 24 giờ trước giờ bắt đầu, toàn hệ thống. \\
08/10/2026 19:58 & \textbf{Q14.} Chuyển người từ waitlist. Khi có chỗ trống, nhóm dự kiến chuyển ngay người đầu waitlist sang danh sách đăng ký rồi gửi thông báo. Nếu gửi thông báo thất bại, người đó vẫn được xem là đã đăng ký chính thức, hay cần thêm trạng thái "chờ xác nhận" ạ? & Khi có chỗ trống: Hệ thống tự động chuyển người đầu tiên trong waitlist vào danh sách đăng ký. Gửi notification cho sinh viên đó. Không cần thêm trạng thái "chờ xác nhận". Nếu notification gửi thất bại, sinh viên vẫn được xem là đã đăng ký chính thức. Notification chỉ là cơ chế thông báo, không phải điều kiện để hoàn tất việc chuyển trạng thái đăng ký. & Tự chuyển người đầu waitlist và gửi notification; thông báo lỗi vẫn tính đã đăng ký; không có trạng thái chờ xác nhận. \\
08/10/2026 19:58 & \textbf{Q15.} Đăng ký trùng. Nhóm dự kiến chặn một sinh viên có nhiều đăng ký còn hiệu lực cho cùng một sự kiện. Nếu sinh viên đã hủy thì có được đăng ký lại không ạ? Nếu sự kiện đã đầy thì người đó vào cuối waitlist có hợp lý không ạ? & Mình đồng ý với nguyên tắc mỗi sinh viên chỉ có một đăng ký còn hiệu lực cho một sự kiện. Nếu sinh viên đã hủy đăng ký thì được phép đăng ký lại, miễn là sự kiện vẫn còn chỗ hoặc còn nhận waitlist. Nếu sự kiện đã đầy thì sinh viên có thể vào cuối waitlist. & Một đăng ký còn hiệu lực/sinh viên/sự kiện; hủy rồi được đăng ký lại, vào cuối waitlist nếu đầy. \\
08/10/2026 19:58 & \textbf{Q19.} Sửa hoặc hủy sự kiện sau khi công bố. Sau khi sự kiện đã được duyệt và công bố, nếu Organizer sửa các thông tin quan trọng (thời gian, địa điểm, số chỗ) thì sự kiện có cần được duyệt lại không ạ? Organizer có được hủy sự kiện đã công bố không; nếu có thì các đăng ký và waitlist hiện có nên xử lý như thế nào? & Organizer được phép chỉnh sửa hoặc hủy sự kiện đã công bố. Trong MVP, không bắt buộc phải xây dựng một quy trình duyệt lại riêng cho mọi thay đổi. Tuy nhiên, với các thay đổi quan trọng như thời gian hoặc địa điểm, hệ thống cần thông báo cho những sinh viên đã đăng ký để họ biết và có thể hủy đăng ký nếu thay đổi không còn phù hợp. Nếu Organizer hủy sự kiện: Sự kiện được chuyển sang trạng thái Cancelled. Các đăng ký hiện tại được hủy theo sự kiện. Waitlist cũng được đóng/hủy. Các sinh viên liên quan nhận notification. & Yêu cầu mới: sửa/hủy sự kiện đã công bố và thông báo người đã đăng ký. Ảnh hưởng A27, A29 và cơ chế notification. \\
08/10/2026 19:58 & \textbf{Q20.} Quyền xem dữ liệu đăng ký và điểm danh. Nhóm dự kiến: Organizer chỉ xem dữ liệu các sự kiện do mình quản lý; Approver/Admin xem được dữ liệu các sự kiện thuộc những đơn vị tổ chức (Khoa, CLB/Đội/Nhóm, phòng ban) mà mình được giao duyệt; sinh viên chỉ xem đăng ký của chính mình. Cách phân quyền này có phù hợp không ạ? & Đồng ý & Phân quyền xem dữ liệu: Organizer xem sự kiện của mình; Admin xem đơn vị được giao; sinh viên xem đăng ký của chính mình. \\
08/10/2026 19:58 & \textbf{Q16.} Milestone 2 và 3. Hiện tại theo như lịch giảng dạy của thầy thì không có 2 nội dung này. Vậy liệu ta có triển khai báo cáo 2 nội dung này không ạ. Nếu có, mong thầy cho nhóm xin ngày dự kiến và đầu ra bắt buộc của Milestone 2 và Milestone 3 được không ạ? Nhóm cần thông tin này để chốt activity network, Gantt và phân bổ nguồn lực. & Không có milestone 2 \& 3, chỉ có 1 lần progress check trên plan & Không có Milestone 2 và 3, chỉ một Progress Check; CP2 là mốc nội bộ của nhóm. \\
08/10/2026 19:58 & \textbf{Q17.} Change event. Mong thầy cho nhóm biết change event dự kiến diễn ra vào khoảng thời gian nào và dưới hình thức nào (yêu cầu mới trên lớp, email hay thay đổi kịch bản) được không ạ? Nhóm có cần dành sẵn dự phòng (reserve) trong Baseline V1.0 cho change event không ạ? & Change Event sẽ được thầy đưa ra trong quá trình thực hiện project theo yêu cầu của môn học; nhóm không cần biết trước nội dung cụ thể. Nhóm nên xem đây là một phần của uncertainty/risk trong project và có thể dành một mức contingency hợp lý trong kế hoạch, nhưng không cần tạo trước một activity cụ thể với tên "Change Event" hoặc giả định trước nội dung thay đổi. & Change event không báo trước, coi là rủi ro và để contingency; không tạo activity tên Change Event. \\
25/10/2026 07:32 & \textbf{Q21.} Phạm vi demo tại Progress Check 09/11. Theo lịch nhóm đã lập sau khi cân đối thời gian của các thành viên, đến ngày 09/11 nhóm sẽ có: đăng nhập và phân quyền theo vai trò; tạo, duyệt và công bố sự kiện; đăng ký và hủy đăng ký theo hạn 24 giờ, có giới hạn số chỗ; waitlist và tự động chuyển người đầu waitlist đang hoàn thiện (dự kiến xong 10/11); bản thử nghiệm quét QR. Điểm danh QR hoàn chỉnh (có nhập mã hoặc xác nhận thủ công dự phòng) dự kiến xong 23--24/11; toàn bộ workflow gồm báo cáo và xuất Excel/CSV được nghiệm thu trong tuần 01--05/12. Tại Progress Check, nhóm cũng sẽ báo cáo actual effort, tiến độ so với kế hoạch và cập nhật risks/issues. Thầy có đồng ý phạm vi trên là đủ cho Progress Check không ạ? Nếu thầy cần thấy điểm danh chạy được tại Progress Check, nhóm sẽ sắp xếp lại thứ tự công việc hoặc cắt bớt yêu cầu mức Should, Could theo thứ tự đã thống nhất (không bỏ workflow cốt lõi), rồi báo lại thầy lịch điều chỉnh. & Mình đồng ý với phạm vi nhóm đề xuất. Tại Progress Check, nhóm chưa bắt buộc phải hoàn thiện toàn bộ workflow, bao gồm QR check-in và báo cáo/Excel/CSV. Tuy nhiên, các chức năng đã hoàn thành cần có bản demo chạy được, không chỉ trình bày giao diện hoặc mô tả thiết kế. Đối với những chức năng đang phát triển, nhóm cần thể hiện rõ tiến độ và kế hoạch hoàn thành. & Giữ lịch, không đổi thứ tự công việc. Tiêu chí CP4: chức năng đã xong phải demo chạy được; chức năng đang làm (waitlist, QR) trình bày tiến độ và kế hoạch hoàn thành. REQ-10, REQ-11 kiểm ở CP5. \\
25/10/2026 07:32 & \textbf{Q22.} Kênh thông báo tối thiểu. Theo trả lời Q14 và Q19 của thầy, hệ thống phải thông báo cho sinh viên khi được chuyển từ waitlist sang danh sách đăng ký, và khi sự kiện đổi thời gian, địa điểm hoặc bị hủy. Nhóm đề xuất thông báo trong ứng dụng là mức bắt buộc của MVP; gửi email là mức Should, chỉ làm nếu còn thời gian. Đề xuất này có phù hợp với thầy không ạ? & Mình đồng ý với đề xuất của nhóm. & Thông báo trong ứng dụng là Must, email là Should (REQ-09, REQ-21); đóng giả định còn mở về kênh thông báo. \\
\end{longtable}
\endgroup
